% ============================================================
%  CHAPITRE 3 — CONCEPTION ET ARCHITECTURE
% ============================================================
\chapter{Conception et Architecture}

\section{Architecture globale du système}

\insureflow{} repose sur une architecture orientée événements composée
de microservices spécialisés communiquant de manière asynchrone via RabbitMQ.
L'ensemble du système est conteneurisé via Docker Compose et déployé localement
sur un MacBook Pro M2 Pro (16 Go RAM) avec accélération GPU déportée via Tailscale
vers un PC équipé d'une RTX 4060 Ti.

\begin{figure}[H]
\centering
\begin{tikzpicture}[x=1cm, y=1cm,
  bbox/.style={rectangle, rounded corners=5pt, draw=primary,
              fill=secondary, minimum width=3cm, minimum height=0.9cm,
              align=center, font=\small\bfseries},
  agent/.style={rectangle, rounded corners=5pt, draw=accent,
                fill=accent!10, minimum width=3.2cm, minimum height=0.9cm,
                align=center, font=\small\bfseries},
  infra/.style={rectangle, rounded corners=5pt, draw=OliveGreen,
                fill=OliveGreen!10, minimum width=3cm, minimum height=0.9cm,
                align=center, font=\small\bfseries},
  arr/.style={-{Stealth[length=5pt]}, thick, color=primary},
]
  \node[bbox] (angular) at (0, 0)    {Angular 17\\Frontend};
  \node[bbox] (postman) at (4, 0)    {Postman /\\API Client};

  \node[bbox, fill=primary!15, draw=primary, minimum width=8.5cm]
    (api) at (2, -1.8) {Spring Boot API — Port 8080};

  \node[infra] (kc) at (8.5, -1.8) {Keycloak 24\\Port 8180};

  \node[infra, minimum width=8.5cm] (rmq) at (2, -3.6)
    {RabbitMQ — Exchange fan-out};

  \node[agent] (router)    at (-1.5, -5.4) {RouterAgent};
  \node[agent] (validator) at ( 2.0, -5.4) {ValidatorAgent};
  \node[agent] (estimator) at ( 5.5, -5.4) {EstimatorAgent};
  \node[agent] (fraud)     at ( 9.0, -5.4) {FraudAgent};

  \node[infra] (pg)      at (-1.5, -7.2) {PostgreSQL\\+ pgvector};
  \node[infra] (ollama)  at ( 2.0, -7.2) {Ollama\\llama3.1 + vision};
  \node[infra] (serp)    at ( 5.5, -7.2) {SerpAPI\\Prix marché};
  \node[infra] (cloud)   at ( 9.0, -7.2) {Cloudinary\\Photos};

  \draw[arr] (angular) -- (api);
  \draw[arr] (postman) -- (api);
  \draw[arr] (api) -- (kc) node[midway, above, font=\tiny] {JWT};
  \draw[arr] (api) -- (rmq);
  \draw[arr] (rmq) -- (router);
  \draw[arr] (rmq) -- (validator);
  \draw[arr] (rmq) -- (estimator);
  \draw[arr] (rmq) -- (fraud);
  \draw[arr] (router)    -- (pg);
  \draw[arr] (validator) -- (ollama);
  \draw[arr] (estimator) -- (serp);
  \draw[arr] (fraud)     -- (cloud);
  \draw[arr] (validator) -- (pg);
  \draw[arr] (estimator) -- (ollama);
\end{tikzpicture}
\caption{Architecture globale du système \insureflow{}}
\label{fig:archi_globale}
\end{figure}

\section{Justification des choix technologiques}

\subsection{Pourquoi une architecture multi-agents ?}

Le traitement d'un sinistre implique quatre responsabilités métier distinctes et
indépendantes : classification, validation contractuelle, estimation financière et
détection de fraude. Une architecture monolithique forcerait ces responsabilités
à partager le même contexte d'exécution, rendant le système rigide et difficile
à faire évoluer.

L'architecture multi-agents permet au contraire de :
\begin{itemize}
  \item \textbf{Isoler les défaillances} : la panne d'un agent n'affecte pas les autres.
  \item \textbf{Exécuter en parallèle} : ValidatorAgent et EstimatorAgent s'exécutent
        simultanément sur le même sinistre, réduisant le temps total de traitement.
  \item \textbf{Spécialiser les modèles} : chaque agent peut utiliser le prompt,
        le modèle et les outils les mieux adaptés à sa tâche.
\end{itemize}

\subsection{Pourquoi RabbitMQ et non REST synchrone ?}

Un appel REST synchrone entre agents crée un couplage fort : si l'EstimatorAgent
met 2 minutes à analyser des photos, l'appelant est bloqué pendant toute cette durée.
RabbitMQ résout ce problème par découplage temporel.

\begin{table}[H]
\centering
\caption{REST synchrone vs RabbitMQ asynchrone}
\begin{tabularx}{\textwidth}{>{\bfseries}lXX}
\toprule
Critère          & REST synchrone & RabbitMQ asynchrone \\
\midrule
Couplage         & Fort (appelant attend) & Faible (pub/sub) \\
Tolérance pannes & Nulle (cascade failure) & Haute (messages persistés) \\
Parallélisme     & Séquentiel & Natif (fan-out exchange) \\
Scalabilité      & Limitée & Horizontale \\
Latence client   & Bloquante & Non-bloquante \\
\bottomrule
\end{tabularx}
\label{tab:rest_vs_rabbit}
\end{table}

\subsection{Pourquoi RAG et pgvector ?}

La vérification de couverture contractuelle ne peut pas être résolue par un LLM
seul : les contrats sont des documents spécifiques à chaque compagnie, non présents
dans les données d'entraînement. La technique RAG permet de :
\begin{itemize}
  \item Vectoriser les clauses contractuelles une seule fois à l'ingestion.
  \item Retrouver en temps réel les clauses pertinentes par similarité cosinus
        via pgvector.
  \item Injecter ces clauses dans le contexte du LLM pour une décision fondée
        sur le contrat réel du client.
\end{itemize}

\subsection{Pourquoi Ollama et LLMs locaux ?}

L'utilisation de modèles locaux via Ollama présente trois avantages décisifs
dans le contexte de l'assurance :
\begin{itemize}
  \item \textbf{Confidentialité} : les données des sinistres ne quittent jamais
        l'infrastructure locale.
  \item \textbf{Coût} : absence de coût par token, critique pour un système
        traitant potentiellement des milliers de sinistres.
  \item \textbf{Contrôle} : possibilité de fine-tuner les modèles sur des données
        métier spécifiques à l'assurance tunisienne.
\end{itemize}

\subsection{Pourquoi Keycloak ?}

Keycloak 24 fournit une solution IAM complète en open source, couvrant
l'authentification OpenID Connect, la gestion des rôles (CLIENT, ADMIN),
la génération de tokens JWT signés RS256 et la gestion des attributs utilisateur
(CIN). Son intégration native avec Spring Boot via Spring Security OAuth2 Resource
Server simplifie considérablement la sécurisation de l'API.

\section{Architecture multi-agents — Conception détaillée}

\subsection{Vue d'ensemble du pipeline}

Chaque sinistre suit un pipeline en cinq étapes orchestrées par RabbitMQ :

\begin{figure}[H]
\centering
\begin{tikzpicture}[x=1cm, y=1cm,
  pstep/.style={rectangle, rounded corners=5pt, draw=primary,
                fill=secondary, minimum width=2.8cm, minimum height=1cm,
                align=center, font=\small\bfseries},
  arr/.style={-{Stealth[length=6pt]}, thick, color=primary},
]
  \node[pstep] (sub)   at (0,  0) {Soumission\\client};
  \node[pstep] (route) at (3,  0) {RouterAgent\\classification};
  \node[pstep] (val)   at (6,  0) {ValidatorAgent\\couverture RAG};
  \node[pstep] (est)   at (9,  0) {EstimatorAgent\\coût + vision};
  \node[pstep] (fraud) at (12, 0) {FraudAgent\\anomalies};

  \node[rectangle, rounded corners=5pt, draw=OliveGreen, fill=OliveGreen!10,
        minimum width=2.8cm, minimum height=1cm, align=center,
        font=\small\bfseries] (dec) at (6, -2.5)
    {DecisionMatrix\\Révision humaine};

  \draw[arr] (sub)   -- (route);
  \draw[arr] (route) -- (val);
  \draw[arr] (val)   -- (est);
  \draw[arr] (est)   -- (fraud);
  \draw[arr] (fraud.south) -- ++(0,-0.8) -- (dec.east);
  \draw[arr] (val.south)   -- ++(0,-0.5)
    node[right, font=\tiny, color=accent] {si non couvert} -- (dec.north);
\end{tikzpicture}
\caption{Pipeline de traitement d'un sinistre dans \insureflow{}}
\label{fig:pipeline}
\end{figure}

\subsection{RouterAgent}

Le RouterAgent est le premier agent du pipeline. Il reçoit la description textuelle
du sinistre via RabbitMQ et utilise \textbf{llama3.1:8b} pour classifier le sinistre
en l'un des types suivants :

\begin{itemize}
  \item \texttt{VEHICLE\_DAMAGE} — dommages sur véhicule automobile
  \item \texttt{PROPERTY\_DAMAGE} — dommages immobiliers (incendie, inondation)
  \item \texttt{THEFT} — vol de véhicule ou de biens
  \item \texttt{HEALTH} — frais médicaux
  \item \texttt{NATURAL\_DISASTER} — catastrophes naturelles
\end{itemize}

Le résultat (type + niveau de confiance) est persisté en base de données et
un événement est publié vers les queues ValidatorAgent et EstimatorAgent.

\subsection{ValidatorAgent}

Le ValidatorAgent vérifie si le sinistre est couvert par le contrat d'assurance
du client. Il implémente la technique RAG en trois étapes :

\begin{enumerate}
  \item \textbf{Recherche vectorielle} : la description du sinistre est vectorisée
        et comparée aux chunks du contrat stockés dans pgvector par similarité cosinus.
  \item \textbf{Extraction des clauses} : les 3 à 5 clauses les plus pertinentes
        sont extraites et injectées dans le contexte du LLM.
  \item \textbf{Décision LLM} : llama3.1:8b analyse la description du sinistre
        au regard des clauses et retourne un booléen \texttt{covered} avec justification.
\end{enumerate}

Si le sinistre n'est pas couvert, le pipeline est court-circuité et le dossier
est directement envoyé en révision humaine avec statut \texttt{NOT\_COVERED}.

\subsection{EstimatorAgent}

L'EstimatorAgent est l'agent le plus complexe du système. Il opère en deux phases :

\paragraph{Phase 1 — Analyse visuelle}
Les photos du sinistre sont téléchargées depuis Cloudinary, encodées en base64
et soumises à \textbf{llama3.2-vision} via Ollama. Le modèle identifie chaque
élément endommagé et lui assigne une sévérité parmi :
\texttt{MINOR}, \texttt{MODERATE}, \texttt{SEVERE}, \texttt{TOTAL\_LOSS}.

\paragraph{Phase 2 — Recherche de prix}
Pour chaque élément endommagé, le service de pricing effectue jusqu'à 5 requêtes
SerpAPI ciblées en français et en anglais, géolocalisées sur la Tunisie.
Les prix extraits sont convertis en TND si nécessaire, puis agrégés en une fourchette
min/max. En cas d'échec total de SerpAPI, un fallback LLM est déclenché,
marquant l'estimation comme \textit{non fiable}.

\subsection{FraudAgent}

Le FraudAgent compare le montant déclaré par le client avec l'estimation calculée
par le système et détecte les anomalies suivantes :

\begin{table}[H]
\centering
\caption{Types d'anomalies détectés par le FraudAgent}
\begin{tabularx}{\textwidth}{>{\bfseries}lXl}
\toprule
Anomalie               & Condition de déclenchement & Score fraude \\
\midrule
PRICE\_INFLATION       & Montant client supérieur à l'estimation système de plus de 50\% & 0.70--0.95 \\
DOCUMENT\_INCONSISTENCY & Incohérence entre description et photos analysées & 0.60--0.85 \\
NONE                   & Aucune anomalie détectée & 0.00--0.10 \\
\bottomrule
\end{tabularx}
\label{tab:fraud_types}
\end{table}

\subsection{OrchestratorConsumer et DecisionMatrix}

L'OrchestratorConsumer centralise les résultats de tous les agents et applique
la DecisionMatrix pour déterminer le statut final du sinistre :

\begin{itemize}
  \item \textbf{Règle 1} — Sinistre non couvert $\rightarrow$ \texttt{REJECTED}.
  \item \textbf{Règle 2} — Score de fraude $\geq$ 0.7 $\rightarrow$ \texttt{PENDING\_REVIEW}.
  \item \textbf{Règle 3} — Sévérité \texttt{TOTAL\_LOSS} $\rightarrow$ \texttt{PENDING\_REVIEW}.
  \item \textbf{Règle 4} — Tous les autres cas $\rightarrow$ \texttt{PENDING\_REVIEW}.
\end{itemize}

\begin{infobox}[Politique de non-approbation automatique]
Aucun sinistre n'est jamais approuvé automatiquement par le système.
Toute décision d'indemnisation requiert une validation explicite d'un administrateur humain.
\end{infobox}

\section{Modèle de données}

\subsection{Diagramme de classes principal}

\begin{figure}[H]
\centering
\begin{tikzpicture}[x=1cm, y=1cm,
  cls/.style={rectangle, draw=primary, fill=white,
              rounded corners=3pt, font=\small, minimum width=4.5cm},
  assoc/.style={thick, color=darkgray},
]
  \node[cls] (client) at (0, 0) {
    \begin{tabular}{l}
      \textbf{\underline{Client}}\\
      \hline
      id : UUID\\
      nom : String\\
      prenom : String\\
      email : String\\
      nationalId : String\\
      keycloakId : String\\
    \end{tabular}
  };

  \node[cls] (policy) at (7, 0) {
    \begin{tabular}{l}
      \textbf{\underline{Policy}}\\
      \hline
      id : UUID\\
      policyNumber : String\\
      type : PolicyType\\
      coverageLimit : Decimal\\
      deductible : Decimal\\
      startDate : Date\\
      endDate : Date\\
    \end{tabular}
  };

  \node[cls] (claim) at (3.5, -5.5) {
    \begin{tabular}{l}
      \textbf{\underline{Claim}}\\
      \hline
      id : UUID\\
      type : ClaimType\\
      status : ClaimStatus\\
      description : String\\
      photoUrls : List\\
      clientEstimatedCost : Decimal\\
      estimatedCost : Decimal\\
      estimatorResult : JSON\\
      fraudResult : JSON\\
      createdAt : DateTime\\
    \end{tabular}
  };

  \node[cls] (doc) at (-0.5, -5.5) {
    \begin{tabular}{l}
      \textbf{\underline{PolicyDocument}}\\
      \hline
      id : UUID\\
      policyId : UUID\\
      content : String\\
      embedding : vector(384)\\
    \end{tabular}
  };

  \draw[assoc] (client.east) -- node[midway,above,font=\tiny]{1\hspace{1cm}*}
    (policy.west);
  \draw[assoc] (client.south) -- node[midway,left,font=\tiny]{1}
    (claim.north west) node[near end,left,font=\tiny]{*};
  \draw[assoc] (policy.south) -- node[midway,right,font=\tiny]{1}
    (claim.north east) node[near end,right,font=\tiny]{*};
  \draw[assoc] (policy.south west) -- node[midway,left,font=\tiny]{1}
    (doc.north) node[near end,left,font=\tiny]{*};
\end{tikzpicture}
\caption{Diagramme de classes principal — \insureflow{}}
\label{fig:classes}
\end{figure}

\section{Diagrammes de séquence}

\subsection{Séquence globale de traitement d'un sinistre}

\begin{figure}[H]
\centering
\resizebox{\textwidth}{!}{%
\begin{tikzpicture}[x=1cm, y=1cm,
  life/.style={draw=primary, thick},
  msg/.style={-{Stealth[length=5pt]}, thick, color=darkgray},
  ret/.style={-{Stealth[length=5pt]}, dashed, thick, color=gray},
  pbox/.style={rectangle, draw=primary, fill=secondary,
              minimum width=2.2cm, minimum height=0.7cm,
              align=center, font=\tiny\bfseries},
]
  \node[pbox] at (0,    0) {Client};
  \node[pbox] at (2.5,  0) {API};
  \node[pbox] at (5,    0) {RabbitMQ};
  \node[pbox] at (7.5,  0) {Router};
  \node[pbox] at (10,   0) {Validator};
  \node[pbox] at (12.5, 0) {Estimator};
  \node[pbox] at (15,   0) {FraudAgent};
  \node[pbox] at (17.5, 0) {Orchestrator};

  \foreach \x in {0,2.5,5,7.5,10,12.5,15,17.5}{
    \draw[life] (\x,-0.35) -- (\x,-11);
  }

  \draw[msg] (0,-0.9)   -- (2.5,-0.9)
    node[midway,above,font=\tiny]{POST /claims};
  \draw[msg] (2.5,-1.4) -- (5,-1.4)
    node[midway,above,font=\tiny]{publish(submitted)};
  \draw[ret] (2.5,-1.9) -- (0,-1.9)
    node[midway,above,font=\tiny]{202 Accepted};

  \draw[msg] (5,-2.5)   -- (7.5,-2.5)
    node[midway,above,font=\tiny]{claim.submitted};
  \draw[msg] (7.5,-3.2) -- (5,-3.2)
    node[midway,above,font=\tiny]{publish(routed)};

  \draw[msg] (5,-3.9)   -- (10,-3.9)
    node[midway,above,font=\tiny]{validate};
  \draw[msg] (5,-4.2)   -- (12.5,-4.2)
    node[midway,above,font=\tiny]{estimate};

  \draw[msg] (10,-5.2)  -- (5,-5.2)
    node[midway,above,font=\tiny]{publish(validated)};
  \draw[msg] (12.5,-6.2) -- (5,-6.2)
    node[midway,above,font=\tiny]{publish(valued)};

  \draw[msg] (5,-7.0)   -- (15,-7.0)
    node[midway,above,font=\tiny]{fraud check};
  \draw[msg] (15,-7.8)  -- (5,-7.8)
    node[midway,above,font=\tiny]{publish(decision)};

  \draw[msg] (5,-8.6)   -- (17.5,-8.6)
    node[midway,above,font=\tiny]{claim.decision};
  \draw[msg] (17.5,-9.4) -- (2.5,-9.4)
    node[midway,above,font=\tiny]{updateStatus(PENDING\_REVIEW)};

  \draw[msg] (0,-10.2)  -- (2.5,-10.2)
    node[midway,above,font=\tiny]{GET /claims/\{id\}};
  \draw[ret] (2.5,-10.7) -- (0,-10.7)
    node[midway,above,font=\tiny]{200 + résultats agents};
\end{tikzpicture}
}
\caption{Diagramme de séquence — Traitement global d'un sinistre}
\label{fig:seq_global}
\end{figure}

\subsection{Séquence d'authentification Keycloak}

\begin{figure}[H]
\centering
\begin{tikzpicture}[x=1cm, y=1cm,
  life/.style={draw=primary, thick},
  msg/.style={-{Stealth[length=5pt]}, thick, color=darkgray},
  ret/.style={-{Stealth[length=5pt]}, dashed, thick, color=gray},
  pbox/.style={rectangle, draw=primary, fill=secondary,
              minimum width=2.5cm, minimum height=0.7cm,
              align=center, font=\tiny\bfseries},
]
  \node[pbox] at (0,  0) {Client};
  \node[pbox] at (4,  0) {Angular};
  \node[pbox] at (8,  0) {Keycloak};
  \node[pbox] at (12, 0) {Spring Boot API};

  \foreach \x in {0,4,8,12}{
    \draw[life] (\x,-0.35) -- (\x,-8.5);
  }

  \draw[msg] (0,-0.9) -- (4,-0.9)
    node[midway,above,font=\tiny]{login(user, pwd)};
  \draw[msg] (4,-1.6) -- (8,-1.6)
    node[midway,above,font=\tiny]{OIDC auth request};
  \draw[ret] (8,-2.3) -- (4,-2.3)
    node[midway,above,font=\tiny]{JWT access\_token};
  \draw[ret] (4,-3.0) -- (0,-3.0)
    node[midway,above,font=\tiny]{token stocké};

  \draw[msg] (0,-3.8)  -- (4,-3.8)
    node[midway,above,font=\tiny]{action utilisateur};
  \draw[msg] (4,-4.5)  -- (12,-4.5)
    node[midway,above,font=\tiny]{GET /api + Bearer JWT};
  \draw[msg] (12,-5.3) -- (8,-5.3)
    node[midway,above,font=\tiny]{valider JWT (RS256)};
  \draw[ret] (8,-6.0)  -- (12,-6.0)
    node[midway,above,font=\tiny]{claims (rôle, CIN)};
  \draw[ret] (12,-6.7) -- (4,-6.7)
    node[midway,above,font=\tiny]{200 + données};
  \draw[ret] (4,-7.4)  -- (0,-7.4)
    node[midway,above,font=\tiny]{affichage};
\end{tikzpicture}
\caption{Diagramme de séquence — Authentification via Keycloak}
\label{fig:seq_auth}
\end{figure}

\section{Infrastructure technique}

\subsection{Stack technique complet}

\begin{table}[H]
\centering
\caption{Stack technique d'\insureflow{}}
\begin{tabularx}{\textwidth}{>{\bfseries}llX}
\toprule
Couche           & Technologie          & Rôle \\
\midrule
Backend          & Spring Boot 3.3      & API REST, orchestration agents \\
Messaging        & RabbitMQ 3.13        & Communication asynchrone inter-agents \\
Auth             & Keycloak 24.0        & IAM, OIDC, JWT RS256 \\
Base de données  & PostgreSQL 16 + pgvector & Persistance + recherche vectorielle \\
LLM local        & Ollama + llama3.1:8b & Routing, validation, fraude \\
Vision           & llama3.2-vision      & Analyse visuelle des photos de sinistres \\
Prix marché      & SerpAPI              & Recherche de prix en temps réel \\
Stockage photos  & Cloudinary           & Upload et CDN des photos \\
Frontend         & Angular 17           & Interface client et panel administrateur \\
Conteneurisation & Docker Compose       & Orchestration locale des services \\
VPN              & Tailscale            & Déport GPU RTX 4060 Ti vers Mac M2 Pro \\
\bottomrule
\end{tabularx}
\label{tab:stack}
\end{table}

\subsection{Configuration RabbitMQ — Exchange fan-out}

\insureflow{} utilise un exchange de type \textbf{fan-out} pour diffuser les
événements à plusieurs agents simultanément. Chaque agent possède sa propre queue
dédiée :

\begin{itemize}
  \item \texttt{claim.router.queue} — consommée par RouterAgent
  \item \texttt{claim.estimated.queue} — consommée par ValidatorAgent et EstimatorAgent
  \item \texttt{claim.fraud.queue} — consommée par FraudAgent
  \item \texttt{claim.decision.queue} — consommée par OrchestratorConsumer
\end{itemize}

\subsection{RAG et base vectorielle (pgvector)}

Les contrats d'assurance sont découpés en chunks de 512 tokens, vectorisés par
le modèle d'embeddings \texttt{all-MiniLM-L6-v2} via LangChain4j et stockés
dans PostgreSQL avec l'extension pgvector. La recherche de clauses pertinentes
utilise la distance cosinus avec un seuil de similarité de 0.7.

\subsection{Analyse visuelle — llama3.2-vision via Ollama}

Les photos des sinistres sont téléchargées depuis Cloudinary, encodées en base64
et soumises à llama3.2-vision avec un prompt structuré imposant un format JSON
strict. Le modèle identifie chaque élément endommagé, lui assigne une sévérité
et produit un raisonnement en français. En l'absence de photos, le système
bascule automatiquement vers une analyse textuelle par llama3.1:8b.

\newpage    